\section{Why How What}

\begin{quotex}
Howbeit we speak wisdom among the perfect: yet not the wisdom of this world, neither of the princes of this world that come to nought; But we speak the wisdom of God in a mystery, a wisdom which is hidden, which God ordained before the world, unto our glory:
\flright{\textsc{1 Corinthians 2:6-7}}
\end{quotex}


I listened to a lecture on TED\footnote{\url{https://www.youtube.com/watch?v=sioZd3AxmnE}} that claimed that an enterprise needs to answer the questions ``\textbf{Why? How? What}?'' in that order. We have answered ``what'' in many different ways, but apparently not the why. So this is an attempt to answer the ``why''.

\paragraph{Fundamental points}


We were motivated at the very beginning by \textbf{Rene Guenon}'s \emph{Crisis of the Modern World}. It makes two fundamental points:


\begin{enumerate}
\item The crisis is the result of the loss of tradition in the west.

\item Tradition will be restored only via an elite
\end{enumerate}


We take point (1) as stipulated, even while we have identified the many specific ways that tradition has been lost. Point (1) has been written about everywhere, without any consistency. So our two ``whys'' follow, but in inverse order.

\begin{enumerate}


\item Gornahoor aims to identify the qualities of a traditional elite and to offer knowledge, suggestions, and opportunities to develop, or initiate, that elite.

\item Gornahoor aims to make clear which specific traditional principles have been undermined by modernity. Following \textbf{Joseph de Maistre}, we don't expect a revolution in the opposite direction (e.g., following some cataclysmic collapse) but rather the opposite of a revolution.
\end{enumerate}


Hence, under the guise of a general interest web site, the real intent is for the relatively few of a like mind to be drawn together in accordance with the Law of Affinity. Therefore, we don't try to convince anyone. We make the case, as necessary, and rely on men's inner convictions to respond as they see fit.

\paragraph{The Nature of Things}


The conventional view is that things obey certain laws external to them which science then discovers. This is not precisely correct, as we have shown that these so-called laws are really creations in the mind of the scientist. It should be clear then that things act according to their own inner nature.

\paragraph{The Nature of Man}


Man, as a person, has intellect and will. That is his nature and is what distinguishes him from other things in the world. Unlike things which are unconscious of their nature, man has an intellect and can come to know his nature. Unlike things which follow their nature automatically, man has a will, so he can choose to follow his nature or not.

\paragraph{Human Action}


The will follows the intellect. Hence a disordered mind will result in disordered actions. That is why the first task is to describe ``what'' is an ordered mind. But everyone makes that same claim; there is no shortage of people who will tell other people what to think. The one thing necessary, then, is the ``how'': how can one achieve an ordered mind and with conviction, knowing it represents the truth.

\paragraph{One Tradition}


We take it as a given that the Catholic tradition is the proper one for the West as Guenon suggests. In the Crisis, he thought that the elite would arise from that, although he was not sanguine about its prospects. However, whatever its current state may be is not really our concern. There is the Church of Peter and the Church of John\footnote{\url{https://www.meditationsonthetarot.com/the-church-of-john}}; it is the latter we follow, although not in conflict with Peter.

We stand with Augustine that there is only one tradition and it is now called Catholicism. That means it must be understood in that light. That will be the How. \emph{Pace} Guenon, we have shown that there is still a living esoteric tradition. However, if anyone is looking for it here or there, he won't find it. There are no buildings, no ads in the backs of magazines, no permanent structures, no outrageous promises. Movements arise spontaneously since the Spirit blows where he wills; we are initiated from above. If you blink, you may miss it.

\paragraph{An Opportunity Lost}


Vatican II offered the opportunity to recover tradition. Unfortunately, it created turmoil, although in the long run things may change. Fantasize for a moment if Rene Guenon had been one of the experts at that Council.

Unfortunately, Rene Guenon was a victim of the super correct convert Jacques Maritain. Maritain managed to exclude Guenon from Catholic intellectual circles, to its intellectual impoverishment. Later on, Maritain seemed to make amends when he mentioned Shankara as a philosopher of being, but at that point it was too late to make any impact on his philosophy. Instead, Maritain made a more liberal turn, becoming an influential figure for Vatican II. Meanwhile, Guenon abandoned that project and turned to Islam.

\paragraph{Signs of the Times}


\begin{quotex}
At all times the Church carries the responsibility of reading the signs of the time. 
\flright{\emph{Gaudium et Spes}}
\end{quotex}


How different would the document on the relationship of the Church to the modern world have turned out had Guenon's book of the same title held sway! We have the responsibility to discern the signs of the times.

\paragraph{Ecumenism}


How much spiritual depth would have been added to the declaration on the relation of the Church to non-Christian religions! Guenon was opposed to false ecumenism and syncretism of the Assisi type where all religions gathered together indiscriminately. Instead, Guenon said that unity lies deeper, in their respective symbolic and metaphysical underpinnings. Ecumenism as described by Guenon is what we pursue.

\paragraph{Ressourcement}


Ressourcement is the call to return to the sources. Again, Guenon has shown us the way. As he pointed out, Platonism, properly understood as a metaphysic rather than a philosophical school, Is a nearly complete metaphysic. He was always quick to point out the many correspondences between \textbf{Aristotle} and \textbf{Plato} on the one side and the various Vedantic systems on the other. Therefore, we, too, are motivated by ressourcement.

Plato, however, was not \emph{sui generis}, but was himself part of a larger Tradition\footnote{\url{https://gornahoor.net/?p=1256}}, neither at the beginning not the end. Hence, we are justified to go back to all those sources, including \textbf{Hermes Trismegistus} and \textbf{Pythagoras}.


As Vladimir Solovyov pointed out, Christianity did not create a new metaphysic. Rather, it is a story of historical events. The proper metaphysic, therefore, can be traced back to Hermetism and Neoplatonism. That is an option to take quite seriously.

Updated 5 September 2020


\flrightit{Posted on 2014-09-05 by Cologero}

